\documentclass[12pt,a4paper]{article}
\newcommand{\imagefolder}{../Figure}
%\newcommand{\imagefolder}{C:/Users/97252/OneDrive - ac.sce.ac.il/Lab Doc/Measuring equipment/DMM/Figure}
%\newcommand{\docconfig}{D:/OneDrive - המכללה האקדמית להנדסה ע ש סמי שמעון/Lab Doc/Measuring equipment/DMM/Doc/doc_config}
%\newcommand{\imagefolder}{D:/OneDrive - המכללה האקדמית להנדסה ע ש סמי שמעון/Lab Doc/Measuring equipment/DMM/Figure}
%\newcommand{\docconfig}{C:/Users/97252/OneDrive - ac.sce.ac.il/Lab Doc/Measuring equipment/DMM/Doc/doc_config}

%% XeLaTeX!
% ============================================================ %
% HEBREW support via polyglossia %
% ============================================================ %
%\usepackage{polyglossia}
\defaultfontfeatures{Mapping=tex-text, Scale=MatchLowercase}
\setdefaultlanguage{english}
\setotherlanguage{hebrew}
\newfontfamily\hebrewfont[Script=Hebrew]{David}
% Use \begin{hebrew} block of text \end{hebrew} for paragraphs.
% Use \texthebrew{ } and \textenglish{ } for short texts.
% ============================================================ %
\newcommand{\te}[1]{\textenglish{#1}}


\newcommand{\thh}[1]{\texthebrew{#1}}
%Hyper Link Setup
\hypersetup{
	colorlinks=true,
	linkcolor=blue,
	filecolor=magenta,      
	urlcolor=blue,
}

%% DMM and link to wiki
\newcommand{\dmm}{\href{https://en.wikipedia.org/wiki/Multimeter}{\te{(DMM)}}}
%

\usepackage{float}
\usepackage{polyglossia}
\usepackage{graphicx,wrapfig,lipsum}
\usepackage{xspace}
\usepackage{glossaries}
\usepackage{color}
\usepackage{fontspec}

\usepackage{blindtext}
\usepackage[utf8]{inputenc}
\usepackage{fancyhdr}

\usepackage{booktabs} %for Excel
\usepackage{multirow}
\usepackage{pbox}

\usepackage{hyperref}

\usepackage{siunitx}
\usepackage{comment}
\usepackage{enumitem}
\usepackage{amsmath}

\usepackage{tabularx} % for table adjustment
%
\usepackage[a4paper ,total={6in, 8in}]{geometry}
%\usepackage[showframe=true]{geometry}


\newcommand{\mysetup}{../../../Latex_set_up}
%\newcommand{\mysetup}{C:/Users/97252/OneDrive - ac.sce.ac.il/Lab Doc/Latex_set_up}


\usepackage{float}
\usepackage{polyglossia}
\usepackage{graphicx,wrapfig,lipsum}
\usepackage{xspace}
\usepackage{glossaries}
\usepackage{color}
\usepackage{fontspec}

\usepackage{blindtext}
\usepackage[utf8]{inputenc}
\usepackage{fancyhdr}

\usepackage{booktabs} %for Excel
\usepackage{multirow}
\usepackage{pbox}

\usepackage{hyperref}

\usepackage{siunitx}
\usepackage{comment}
\usepackage{enumitem}
\usepackage{amsmath}

\usepackage{tabularx} % for table adjustment
%
\usepackage[a4paper ,total={6in, 8in}]{geometry}
%\usepackage[showframe=true]{geometry}
% XeLaTeX!
% ============================================================ %
% HEBREW support via polyglossia %
% ============================================================ %
%\usepackage{polyglossia}
\defaultfontfeatures{Mapping=tex-text, Scale=MatchLowercase}
\setdefaultlanguage{english}
\setotherlanguage{hebrew}
\newfontfamily\hebrewfont[Script=Hebrew]{David}
% Use \begin{hebrew} block of text \end{hebrew} for paragraphs.
% Use \texthebrew{ } and \textenglish{ } for short texts.
% ============================================================ %
\newcommand{\te}[1]{\textenglish{#1}}


\newcommand{\thh}[1]{\texthebrew{#1}}
%Hyper Link Setup
\hypersetup{
	colorlinks=true,
	linkcolor=blue,
	filecolor=magenta,      
	urlcolor=blue,
}

%% DMM and link to wiki
\newcommand{\dmm}{\href{https://en.wikipedia.org/wiki/Multimeter}{\te{(DMM)}}}


%\urlstyle{same}


\begin{document}
	\begin{hebrew}
		\title{רב-מודד ספרתי \dmm}
%		\author{עודד פוצ'ינסקי}
		\maketitle
	\end{hebrew}


\begin{hebrew}
	\textbf{רב מודד דיגיטאלי - \te{DMM}}
	\textbf{: }
	הינו מכשיר מדידה רב תכליתי המסוגל למדוד מתח, זרם, והתנגדות. רב המודד מסומן ב DMM
	שפירושם Meter Multi Digital .לרב המודד יש בדרך כלל 3 זוגות הדקים נפרדים, אך לרוב לרב המודד יהיה הדק משותף בשם
	COMMON המסומן בצבע שחור. זוג הדקים ראשון משמש למדידת מתח והתנגדות. 
\end{hebrew}



\begin{hebrew}
\section{סוגי המכשירים במעבדה}
	\begin{enumerate}
		\item 
		\te{GDM - 8341}
				\begin{figure}[h]
					\centering
					\includegraphics[scale=0.045]{\imagefolder/GDM_8341/GDM-8341.png}
				\end{figure}
		\item 
		\te{GDM - 8245}
		\begin{figure}[h]
			\centering
			\includegraphics[scale=0.05]{\imagefolder/GDM_8245/GDM-8245.png}
		\end{figure}
	\end{enumerate}
\end{hebrew}

\pagebreak
	
\begin{hebrew}
	\section{פירוט על כל דגם}
		\subsection{GDM - 8341}
	\begin{figure}[h]
		\centering
		\includegraphics[scale=0.05]{\imagefolder/GDM_8341/GDM-8341.png}
	\end{figure}	
			\begin{itemize}
				\item
				\begin{minipage}{0.2\textwidth}
					\centering
					\includegraphics[scale=0.9]{\imagefolder/GDM_8341/INPUT_ V_8341.png}
				\end{minipage}
				\hfill
				\begin{minipage}{0.7\textwidth}
					חיבור המשמש לכל המדידות חוץ מדידות זרם
					\te{DC/AC}	
				\end{minipage}
				\item
				\begin{minipage}{0.2\textwidth}
					\centering
					\includegraphics[scale=0.9]{\imagefolder/GDM_8341/INPUT_COMP_8341.png}
				\end{minipage}
				\hfill
				\begin{minipage}{0.7\textwidth}
					חיבור האדמה	
				\end{minipage}
				\item
				\begin{minipage}{0.2\textwidth}
					\centering
					\includegraphics[scale=0.9]{\imagefolder/GDM_8341/INPUT_LOW_A_8341.png}
				\end{minipage}
				\hfill
				\begin{minipage}{0.7\textwidth}
					חיבור לזרמים נמוכים של עד ל
					$\SI{0.5}{\ampere}$	
				\end{minipage}
				
				\item
				\begin{minipage}{0.2\textwidth}
					\centering
					\includegraphics[scale=0.9]{\imagefolder/GDM_8341/INPUT_A_8341.png}
				\end{minipage}
				\hfill
				\begin{minipage}{0.7\textwidth}
					חיבור לזרמים  עד ל 
					$\SI{12}{\ampere}$ 
					
					
				\end{minipage}
				\item
				\begin{minipage}{0.2\textwidth}
					\centering
					\includegraphics[scale=0.9]{\imagefolder/GDM_8341/DCV_8341.png}
				\end{minipage}
				\hfill
				\begin{minipage}{0.7\textwidth}
					מודד  \textbf{מתח} DC (מודד מתח ממוצע כאשר מודדים על אות \te{AC})
					
					
				\end{minipage}
				\item
				
				\begin{minipage}{0.2\textwidth}
					\centering
					\includegraphics[scale=0.9]{\imagefolder/GDM_8341/DCI_8341.png}
				\end{minipage}
				\hfill
				\begin{minipage}{0.7\textwidth}
					מודד  \textbf{זרם} DC
				\end{minipage}
				\item
				
				\begin{minipage}{0.2\textwidth}
					\centering
					\includegraphics[scale=0.9]{\imagefolder/GDM_8341/ACV_8341.png}
				\end{minipage}
				\hfill
				\begin{minipage}{0.7\textwidth}
					מודד  \textbf{מתח} AC (מודד מתח \te{RMS})
				\end{minipage}
				\item
				\begin{minipage}{0.2\textwidth}
					\centering
					\includegraphics[scale=0.9]{\imagefolder/GDM_8341/ACI_8341.png}
				\end{minipage}
				\hfill
				\begin{minipage}{0.7\textwidth}
					מודד  \textbf{זרם} AC
				\end{minipage}
				\item
				\begin{minipage}{0.2\textwidth}
					\centering
					\includegraphics[scale=0.9]{\imagefolder/GDM_8341/DCV+ACV_8341.png}
				\end{minipage}
				\hfill
				\begin{minipage}{0.7\textwidth}
					מודד \textbf{מתח}
					\te{DC+AC}
					\\
					כלומר מודד \textbf{מתח} 
					\te{true RMS}	
				\end{minipage}
				\item
				\begin{minipage}{0.35\textwidth}
					\centering
					\includegraphics[scale=0.9]{\imagefolder/GDM_8341/DCI+ACI_8341.png}
				\end{minipage}
				\hfill
				\begin{minipage}{0.55\textwidth}
					מודד \textbf{זרם} 
					\te{DC+AC}
					
				\end{minipage}
				\item
				\begin{minipage}{0.2\textwidth}
					\centering
					\includegraphics[scale=0.9]{\imagefolder/GDM_8341/R_8341.png}
				\end{minipage}
				\hfill
				\begin{minipage}{0.7\textwidth}
					מודד \textbf{התנגדות} \\
					ואם לוחצים לחיצה נוספת מודד \textbf{מוליכות} (בעזרת לחיצה נוספת ניתן לחזור חזרה למדידת התנגדות)
				\end{minipage}
				\item
				\begin{minipage}{0.3\textwidth}
					\centering
					\includegraphics[scale=0.9]{\imagefolder/GDM_8341/R_dB_8341.png}
				\end{minipage}
				\hfill
				\begin{minipage}{0.6\textwidth}
					מודד 
					\textbf{\te{dB}}
				\end{minipage}
				\item
				\begin{minipage}{0.3\textwidth}
					\centering
					\includegraphics[scale=0.9]{\imagefolder/GDM_8341/D_C_8341.png}
				\end{minipage}
				\hfill
				\begin{minipage}{0.6\textwidth}
					בודק \textbf{דיודות} \\
					ואם לוחצים לחיצה נוספת ,אז מודד \textbf{קיבול}  (בעזרת לחיצה נוספת ניתן לחזור חזרה לבדיקת דיודות)
				\end{minipage}
				\item
				\begin{minipage}{0.3\textwidth}
					\centering
					\includegraphics[scale=0.9]{\imagefolder/GDM_8341/dBm_8341.png}
				\end{minipage}
				\hfill
				\begin{minipage}{0.6\textwidth}
					מודד 
					\textbf{\te{dBm}}
				\end{minipage}
				\item
				\begin{minipage}{0.3\textwidth}
					\centering
					\includegraphics[scale=0.9]{\imagefolder/GDM_8341/Hz_P_8341.png}
				\end{minipage}
				\hfill
				\begin{minipage}{0.6\textwidth}
					מודד \textbf{תדר} \\
					ואם לוחצים פעם נוספת אז הוא מודד \textbf{מחזור של אות}
				\end{minipage}
				\item
				\begin{minipage}{0.3\textwidth}
					\centering
					\includegraphics[scale=0.9]{\imagefolder/GDM_8341/TEMP_8341.png}
				\end{minipage}
				\hfill
				\begin{minipage}{0.6\textwidth}
					מודד \textbf{טמפרטורה}
				\end{minipage}
				\item
				\begin{minipage}{0.3\textwidth}
					\centering
					\includegraphics[scale=0.9]{\imagefolder/GDM_8341/2nd_8341.png}
				\end{minipage}
				\hfill
				\begin{minipage}{0.6\textwidth}
					{\huge \textcolor{red}{חסר}}
				\end{minipage}
				\item
				\begin{minipage}{0.3\textwidth}
					\centering
					\includegraphics[scale=0.9]{\imagefolder/GDM_8341/REL_8341.png}
				\end{minipage}
				\hfill
				\begin{minipage}{0.6\textwidth}
					{\huge \textcolor{red}{חסר}}
				\end{minipage}
				\item
				\begin{minipage}{0.3\textwidth}
					\centering
					\includegraphics[scale=0.9]{\imagefolder/GDM_8341/REL1_8341.png}
				\end{minipage}
				\hfill
				\begin{minipage}{0.6\textwidth}
					{\huge \textcolor{red}{חסר}}
				\end{minipage}
				\item
				\begin{minipage}{0.3\textwidth}
					\centering
					\includegraphics[scale=0.9]{\imagefolder/GDM_8341/MX_MN_8341.png}
				\end{minipage}
				\hfill
				\begin{minipage}{0.6\textwidth}
					מודד \textbf{ערך מקסימלי} \\
					ואם לוחצים פעם נוספת אז מודד \textbf{ערך מינימלי} 
				\end{minipage}
				\item
				\begin{minipage}{0.3\textwidth}
					\centering
					\includegraphics[scale=0.9]{\imagefolder/GDM_8341/MATH_8341.png}
				\end{minipage}
				\hfill
				\begin{minipage}{0.6\textwidth}
					{\huge \textcolor{red}{חסר}}
				\end{minipage}
				\item
				\begin{minipage}{0.3\textwidth}
					\centering
					\includegraphics[scale=0.9]{\imagefolder/GDM_8341/HOLD_8341.png}
				\end{minipage}
				\hfill
				\begin{minipage}{0.6\textwidth}
					מפעיל את פעולת ההחזקה
				\end{minipage}
				\item
				\begin{minipage}{0.3\textwidth}
					\centering
					\includegraphics[scale=0.9]{\imagefolder/GDM_8341/COMP_8341.png}
				\end{minipage}
				\hfill
				\begin{minipage}{0.6\textwidth}
					מפעיל את פעולת ההשוואה מדידות
				\end{minipage}
				\item
				\begin{minipage}{0.3\textwidth}
					\centering
					\includegraphics[scale=0.9]{\imagefolder/GDM_8341/TRIG_8341.png}
				\end{minipage}
				\hfill
				\begin{minipage}{0.6\textwidth}
					{\huge \textcolor{red}{חסר}}
				\end{minipage}
				\item
				\begin{minipage}{0.3\textwidth}
					\centering
					\includegraphics[scale=0.9]{\imagefolder/GDM_8341/INT_EXT_8341.png}
				\end{minipage}
				\hfill
				\begin{minipage}{0.6\textwidth}
					{\huge \textcolor{red}{חסר}}
				\end{minipage}
				\item
				\begin{minipage}{0.3\textwidth}
					\centering
					\includegraphics[scale=0.9]{\imagefolder/GDM_8341/MENU_8341.png}
				\end{minipage}
				\hfill
				\begin{minipage}{0.6\textwidth}
					נכנס לתפריט \\
					הגדרות מערכת,הגדרות של מדידות ,וכו'
				\end{minipage}
				\item
				\begin{minipage}{0.3\textwidth}
					\centering
					\includegraphics[scale=0.9]{\imagefolder/GDM_8341/SHIFT_EXIT_8341.png}
				\end{minipage}
				\hfill
				\begin{minipage}{0.6\textwidth}
					כש לוחצים על הכפתור 
					\te{SHIFT}
					מקבלים גישה לכל הפעולות  מדידה המשניות . \\
					
					וכש נמצאים בתפריט הכפתור משמש בתור לחצן יציאה .								
				\end{minipage}
				\item
				\begin{minipage}{0.3\textwidth}
					\centering
					\includegraphics[scale=0.9]{\imagefolder/GDM_8341/AUTO_ENTER_8341.png}
				\end{minipage}
				\hfill
				\begin{minipage}{0.6\textwidth}
					כש משתמשים בכפתור בתור 
					\te{AUTO}
					המכשיר יגדיר אוטמטית את הטווח של הפעולה הנבחרת. \\
					וכש משתמשים בכפתור בתור
					\te{ENTER}
					המכשיר מאשר את הערך הנבחר בתוך התפריט
					
				\end{minipage}
				\item
				\begin{minipage}{0.3\textwidth}
					\centering
					\includegraphics[scale=0.9]{\imagefolder/GDM_8341/ARROW_KEYS_8341.png}
				\end{minipage}
				\hfill
				\begin{minipage}{0.6\textwidth}
					החצים משמשים לכדי לנוע בתפריט ולערוך ערכים. \\
					החצים למעלה ולמטה משמשים להגדרת הטווח בצורה ידנית של מדידות. \\
					החצים ימינה שמאלה הם בנוסף משנים את קצב הרענון בין מהיר ,בינוני ואיטי. 
				\end{minipage}	
			\end{itemize}
		
			\begin{hebrew}
				לפירוט נוסף אפשר לפנות ל:
				\href{https://www.manufakturet.create.aau.dk/digitalAssets/858/858521_gdm8341-gdm-8342-um.pdf}{\te{User Manual}}
			\end{hebrew}
\pagebreak
		\subsection{GDM - 8245}
\begin{figure}[h]
	\centering
	\includegraphics[scale=0.05]{\imagefolder/GDM_8245/GDM-8245.png}
\end{figure}	
\begin{itemize}
	\item
		\begin{minipage}{0.3\textwidth}
			\centering
			\includegraphics[scale=0.5]{\imagefolder/GDM_8245/INPUT_ V_8245.png}
		\end{minipage}
		\hfill
		\begin{minipage}{0.6\textwidth}
			חיבור המשמש לכל המדידות חוץ מדידות זרם
			\te{DC/AC}	
		\end{minipage}
	\item
		\begin{minipage}{0.3\textwidth}
			\centering
			\includegraphics[scale=0.5]{\imagefolder/GDM_8245/INPUT_COMP_8245.png}
		\end{minipage}
		\hfill
		\begin{minipage}{0.6\textwidth}
			חיבור האדמה	
		\end{minipage}
	\item
		\begin{minipage}{0.3\textwidth}
			\centering
			\includegraphics[scale=0.5]{\imagefolder/GDM_8245/INPUT_LOW_A_8245.png}
		\end{minipage}
		\hfill
		\begin{minipage}{0.6\textwidth}
			חיבור לזרמים נמוכים של עד ל
			$\SI{0.5}{\ampere}$	
		\end{minipage}
	\item
		\begin{minipage}{0.3\textwidth}
			\centering
			\includegraphics[scale=0.5]{\imagefolder/GDM_8245/INPUT_A_8245.png}
		\end{minipage}
		\hfill
		\begin{minipage}{0.6\textwidth}
			חיבור לזרמים  עד ל
			$\SI{20}{\ampere}$
		\end{minipage}
	\item
		\begin{minipage}{0.2\textwidth}
			\centering
			\includegraphics[scale=0.9]{\imagefolder/GDM_8245/DCV_8245.png}
		\end{minipage}
		\hfill
		\begin{minipage}{0.7\textwidth}
			מודד  \textbf{מתח} DC (מודד מתח ממוצע כאשר מודדים על אות \te{AC}) \\
			כדי למדוד ב 
			$\si{\milli\volt}$ 
			יש ללחוץ לפני על 
			\te{SHIFT}
			לפני.
		
		\end{minipage}
	\item	
			\begin{minipage}{0.2\textwidth}
				\centering
				\includegraphics[scale=0.9]{\imagefolder/GDM_8245/DCI_8245.png}
			\end{minipage}
			\hfill
			\begin{minipage}{0.7\textwidth}
				מודד  \textbf{זרם}  DC \\	
				כדי למדוד עד ל	$\SI{20}{\ampere}$	יש ללחוץ על		
				\te{SHIFT}
				 לפני.
		
		\end{minipage}
	\item			
			\begin{minipage}{0.2\textwidth}
				\centering
				\includegraphics[scale=0.9]{\imagefolder/GDM_8245/ACV_8245.png}
			\end{minipage}
			\hfill
			\begin{minipage}{0.7\textwidth}
				מודד  \textbf{מתח} AC (מודד מתח \te{RMS}) \\
				כדי למדוד ב 
				$\si{\milli\volt}$ 
				יש ללחוץ לפני על 
				\te{SHIFT}
				לפני.
				
		\end{minipage}
	\item
		\begin{minipage}{0.2\textwidth}
			\centering
			\includegraphics[scale=0.9]{\imagefolder/GDM_8245/ACI_8245.png}
		\end{minipage}
		\hfill
		\begin{minipage}{0.7\textwidth}
			מודד  \textbf{זרם} AC \\
			כדי למדוד עד ל	$\SI{20}{\ampere}$	יש ללחוץ על		
			\te{SHIFT}
			לפני.
		\end{minipage}
	\item
	\begin{minipage}{0.2\textwidth}
		\centering
		\includegraphics[scale=0.9]{\imagefolder/GDM_8245/DCV+ACV_8245.png}
	\end{minipage}
	\hfill
		\begin{minipage}{0.7\textwidth}
			מודד \textbf{מתח}
			\te{DC+AC}
			(כלומר מודד \textbf{מתח}
			
			\te{true RMS}) \\
			כדי למדוד תדירות יש ללחוץ על 
			\te{SHIFT}
		\end{minipage}

	\item
		\begin{minipage}{0.2\textwidth}
			\centering
			\includegraphics[scale=0.9]{\imagefolder/GDM_8245/R_8245.png}
		\end{minipage}
		\hfill
		\begin{minipage}{0.7\textwidth}
			מודד \textbf{התנגדות} 
		\end{minipage}
	\item
		\begin{minipage}{0.2\textwidth}
			\centering
			\includegraphics[scale=0.9]{\imagefolder/GDM_8245/CON_8245.png}
		\end{minipage}
		\hfill
		\begin{minipage}{0.7\textwidth}
			מודד \textbf{מוליכות} \\
			כדי לבדוק כיוון של דיודה יש ללחוץ לפני על
			\te{SHIFT}
		\end{minipage}
	\item
		\begin{minipage}{0.2\textwidth}
			\centering
			\includegraphics[scale=0.9]{\imagefolder/GDM_8245/MX_MN_8245.png}
		\end{minipage}
		\hfill
		\begin{minipage}{0.7\textwidth}
			מודד \textbf{ערך מקסימלי} \\
			ואם לוחצים פעם נוספת אז מודד \textbf{ערך מינימלי} 
		\end{minipage}	
	\item
		\begin{minipage}{0.2\textwidth}
			\centering
			\includegraphics[scale=0.9]{\imagefolder/GDM_8245/C_8245.png}
		\end{minipage}
		\hfill
		\begin{minipage}{0.7\textwidth}
			מודד \textbf{קיבול} \\
			
			 כדי למדוד 
			\textbf{\te{dB}}
			יש ללחוץ לפני על 
			\te{SHIFT}
		\end{minipage}
	\item
		\begin{minipage}{0.2\textwidth}
			\centering
			\includegraphics[scale=0.9]{\imagefolder/GDM_8245/HOLD_8245.png}
		\end{minipage}
		\hfill
		\begin{minipage}{0.7\textwidth}
			מפעיל את פעולת ההחזקה \\
			כדי להפעיל את REL יש ללחוץ לפני על 
			\te{SHIFT}
			
			
		\end{minipage}
	\item
		\begin{minipage}{0.2\textwidth}
			\centering
			\includegraphics[scale=0.9]{\imagefolder/GDM_8245/AUTO_ENTER_8245.png}
		\end{minipage}
		\hfill
		\begin{minipage}{0.7\textwidth}
			כש משתמשים בכפתור בתור 
			\te{AUTO}
			המכשיר יגדיר אוטמטית את הטווח של הפעולה הנבחרת. \\
			וכש משתמשים בכפתור בתור
			\te{MENU}
			המכשיר יפתח את התפריט
			
		\end{minipage}
	\item
		\begin{minipage}{0.2\textwidth}
			\centering
			\includegraphics[scale=0.9]{\imagefolder/GDM_8245/SHIFT_8245.png}
		\end{minipage}
		\hfill
		\begin{minipage}{0.7\textwidth}
			כש לוחצים על הכפתור 
			\te{SHIFT}
			מקבלים גישה לכל הפעולות  מדידה המשניות . \\
								
		\end{minipage}

	\item
		\begin{minipage}{0.2\textwidth}
			\centering
			\includegraphics[scale=0.9]{\imagefolder/GDM_8245/ARROW_KEY_8245.png}
		\end{minipage}
		\hfill
		\begin{minipage}{0.7\textwidth}
			החצים משמשים לכדי לנוע בתפריט ולערוך ערכים. \\
			החצים למעלה ולמטה משמשים להגדרת הטווח בצורה ידנית של מדידות.
		\end{minipage}	
\end{itemize}
	\begin{hebrew}
	לפירוט נוסף אפשר לפנות ל:
	\href{https://datasheet.octopart.com/GDM-8245-Instek-datasheet-11997973.pdf}{\te{User Manual}}
	\end{hebrew}
				
\end{hebrew}			
		

	
	
\end{document}
